\documentclass[11pt]{article}
\newcommand{\DocShort}{Rúbrica del equipo expositor}
\usepackage{fontspec}
\setmainfont{TeX Gyre Pagella}
\setsansfont{TeX Gyre Heros}
\usepackage[spanish,es-noshorthands]{babel}
\decimalpoint
\usepackage{geometry}
\geometry{letterpaper, top=2.4cm, bottom=2.4cm, left=2.7cm, right=2.7cm}
\usepackage[expansion=false]{microtype}
\usepackage{parskip}
\setlength{\parskip}{6pt}
\usepackage{amsmath,amssymb}
\usepackage{booktabs}
\usepackage{array}
\usepackage{longtable}
\usepackage{caption}
\usepackage[dvipsnames,table]{xcolor}
\usepackage{enumitem}
\setlist[itemize]{topsep=2pt, itemsep=2pt, parsep=0pt}
\setlist[enumerate]{topsep=2pt, itemsep=2pt, parsep=0pt}
\usepackage{fancyhdr}
\usepackage[hidelinks,pdfencoding=auto]{hyperref}
\usepackage{titlesec}
\usepackage{tcolorbox}
\tcbuselibrary{breakable}
\usepackage{tabularx}
\usepackage{makecell}
\usepackage{ragged2e}
\usepackage{csquotes}
\usepackage[backend=biber,style=apa,sorting=nyt,sortcites=true,language=spanish]{biblatex}
\DeclareLanguageMapping{spanish}{spanish-apa}
\addbibresource{referencias_presentaciones.bib}

\definecolor{darkblue}{HTML}{174A63}
\definecolor{medblue}{HTML}{2E7896}
\definecolor{lightblue}{HTML}{DCEEF4}
\definecolor{teal}{HTML}{1B7F79}
\definecolor{lightteal}{HTML}{D7EFEC}
\definecolor{lightgrey}{HTML}{F5F7F8}
\definecolor{midgrey}{HTML}{5B6770}
\definecolor{warnyellow}{HTML}{FFF4CC}

\titleformat{\section}
{\large\bfseries\color{darkblue}}
{\thesection}{0.75em}{}[{\color{medblue}\titlerule[0.8pt]}]
\titleformat{\subsection}
{\normalsize\bfseries\color{darkblue}}
{\thesubsection}{0.75em}{}
\titleformat{\subsubsection}
{\normalsize\bfseries\color{medblue}}
{\thesubsubsection}{0.75em}{}

\pagestyle{fancy}
\fancyhf{}
\fancyhead[L]{\small\color{midgrey}ICV-342-T Hidrología}
\fancyhead[R]{\small\color{midgrey}\DocShort}
\fancyfoot[C]{\small\color{midgrey}\thepage}
\setlength{\headheight}{14pt}
\renewcommand{\headrulewidth}{0.4pt}
\renewcommand{\footrulewidth}{0pt}

\rowcolors{2}{lightgrey}{white}
\renewcommand{\arraystretch}{1.18}
\newcolumntype{Y}{>{\RaggedRight\arraybackslash}X}

\newtcolorbox{infobox}{
  breakable,
  colback=lightblue,
  colframe=medblue,
  boxrule=0.8pt,
  arc=3pt,
  left=8pt,
  right=8pt,
  top=7pt,
  bottom=7pt
}

\newtcolorbox{warnbox}{
  breakable,
  colback=warnyellow,
  colframe=orange!70!black,
  boxrule=0.8pt,
  arc=3pt,
  left=8pt,
  right=8pt,
  top=7pt,
  bottom=7pt
}

\newcommand{\doccover}[2]{
  \begin{center}
  {\small\color{midgrey}Pontificia Universidad Católica Madre y Maestra\\Escuela de Ingeniería Civil y Ambiental}

  \vspace{1.0cm}
  {\color{medblue}\rule{\linewidth}{2pt}}
  \vspace{0.3cm}

  {\Large\bfseries\color{darkblue}ICV-342-T Hidrología}\\[0.35em]
  {\LARGE\bfseries\color{darkblue}#1}

  \vspace{0.3cm}
  {\color{medblue}\rule{\linewidth}{2pt}}
  \vspace{0.9cm}
  \end{center}

  \begin{center}
  \begin{tabular}{ll}
  \toprule
  \textbf{Asignatura} & ICV-342-T Hidrología \\
  \textbf{Profesor} & Ricardo Rafael Hernández Moreira, PhD \\
  \textbf{Periodo académico} & Mayo--agosto de 2026 \\
  \textbf{Versión} & #2 \\
  \bottomrule
  \end{tabular}
  \end{center}

  \vspace{0.4cm}
}

\newcommand{\unitbox}[3]{
  \begin{tcolorbox}[
    breakable,
    colback=lightteal,
    colframe=teal,
    boxrule=0.8pt,
    arc=3pt,
    title=\textbf{#1}
  ]
  \textbf{Enfoque de la exposición.} #2

  \vspace{0.5\baselineskip}
  \textbf{Apoyo bibliográfico sugerido.} #3
  \end{tcolorbox}
}

\AtBeginDocument{
  \hypersetup{
    pdftitle={\DocShort},
    pdfauthor={Ricardo Rafael Hernández Moreira},
    pdfsubject={ICV-342-T Hidrología}
  }
}


\begin{document}
\doccover{Rúbrica del equipo expositor}{26 de mayo de 2026}

\begin{infobox}
Este formato se utiliza para calificar la exposición oral de cada equipo. La evidencia Exposición oral vale 5\% en el programa de la asignatura y la parte correspondiente al equipo expositor equivale a 4\% del curso.
\end{infobox}

\begin{warnbox}
La duración esperada de la exposición es de 35 a 40 minutos. Si la exposición se extiende demasiado, afecta la discusión y el cierre docente. Si resulta excesivamente corta, se aplicarán las penalizaciones definidas en este mismo formato.
\end{warnbox}

\section{Distribución interna}
\begin{center}
\begin{tabular}{p{0.44\linewidth} p{0.18\linewidth} p{0.22\linewidth}}
\toprule
\textbf{Componente} & \textbf{Peso en el curso} & \textbf{Observación} \\
\midrule
Desempeño del equipo expositor & 4\% & Se califica con la rúbrica de este documento. \\
Participación de audiencia y revisión breve por sesión & 1\% & Se registra aparte mediante el formulario individual de la audiencia. \\
\bottomrule
\end{tabular}
\end{center}

\section{Rúbrica}
\small
{\setlength{\tabcolsep}{4pt}
\begin{longtable}{@{}p{0.22\linewidth} p{0.17\linewidth} p{0.17\linewidth} p{0.17\linewidth} p{0.17\linewidth}@{}}
\toprule
\textbf{Criterio} & \textbf{Excelente} & \textbf{Satisfactorio} & \textbf{En desarrollo} & \textbf{Insuficiente} \\
\midrule
\endfirsthead
\toprule
\textbf{Criterio} & \textbf{Excelente} & \textbf{Satisfactorio} & \textbf{En desarrollo} & \textbf{Insuficiente} \\
\midrule
\endhead
Cobertura del tema (20\%) & Atiende el núcleo del programa, distingue ideas centrales y no omite piezas necesarias para comprender la unidad. & Cubre la mayor parte del tema, con omisiones menores o poco impacto global. & La cobertura es parcial o desbalanceada. & Omite partes esenciales o desordena el tema de forma seria. \\
Dominio técnico (25\%) & Explica criterios, relaciones y supuestos con seguridad; responde preguntas con fundamento técnico. & Muestra dominio razonable con vacíos menores. & Se observan dudas importantes o respuestas poco precisas. & El equipo no demuestra control suficiente del contenido. \\
Claridad y estructura (15\%) & La secuencia de ideas es nítida, bien jerarquizada y fácil de seguir. & La exposición es entendible, aunque con transiciones o énfasis mejorables. & La organización complica el seguimiento del tema. & La exposición es confusa o impropia para aprendizaje del grupo. \\
Aplicación e interpretación (15\%) & Integra ejemplo, datos, gráfico o situación de ingeniería y extrae conclusiones útiles. & Incluye alguna aplicación pertinente, aunque limitada. & La aplicación aparece débil o poco conectada con el contenido. & No aterriza el contenido en un uso técnico reconocible. \\
Material visual (10\%) & Las diapositivas son legibles, limpias y técnicas; cada imagen o figura usada tiene función clara, se comenta en clase y está citada. & El material cumple, aunque con algunos problemas de legibilidad, síntesis o trazabilidad visual. & El apoyo visual dificulta parcialmente la comprensión o usa recursos gráficos poco justificados. & El material es desordenado, ilegible, improductivo o usa imágenes decorativas o no citadas. \\
Manejo del tiempo y trabajo de equipo (15\%) & El tiempo se controla bien, hay participación equilibrada y transiciones razonables. & El manejo del tiempo o la distribución del trabajo presenta debilidades menores. & Hay desequilibrio notable o problemas claros de tiempo. & La coordinación del equipo afecta seriamente la calidad de la sesión. \\
\bottomrule
\end{longtable}
}
\normalsize

\section{Penalización por duración excesivamente corta}
\begin{warnbox}
Como regla operativa:
\begin{itemize}
\item Si la exposición efectiva dura entre 30 y 34 minutos, se descontará 10\% de la calificación obtenida en la rúbrica del equipo expositor.
\item Si la exposición efectiva dura entre 25 y 29 minutos, se descontará 20\% de la calificación obtenida en la rúbrica del equipo expositor.
\item Si la exposición efectiva dura menos de 25 minutos, se descontará 35\% de la calificación obtenida en la rúbrica del equipo expositor.
\end{itemize}
Terminar antes por pocos segundos no activa penalización; la referencia es la duración efectiva total de la exposición académica del equipo.
\end{warnbox}

\section{Observaciones de uso}
\begin{itemize}
\item Este formato evalúa únicamente al equipo expositor.
\item La participación de la audiencia se califica mediante un instrumento aparte.
\item La puntuación final del equipo debe considerar la penalización por duración si corresponde.
\end{itemize}

\end{document}
